\documentclass[UTF8]{ctexart}

\usepackage[a4paper,margin=2.4cm]{geometry}
\usepackage{hyperref}
\usepackage{xcolor}
\usepackage{enumitem}
\usepackage{booktabs}
\usepackage{amsmath}
\usepackage{listings}

\hypersetup{
  colorlinks=true,
  linkcolor=blue!55!black,
  urlcolor=blue!55!black,
  citecolor=blue!55!black
}

\lstset{
  basicstyle=\ttfamily\small,
  breaklines=true,
  columns=fullflexible,
  frame=single
}

\title{Racer-LLM 轨道尺度、宽度与投影问题排查日志}
\author{Codex 与实验记录}
\date{2026-06-16}

\begin{document}
\maketitle

\section{问题背景}

当前实验目标是使用最简单的 MPC 控制单车跟随读取的 raceline，并尽量保持车辆不出界。
然而在 raw CSV 坐标系下使用赛道半宽 $0.22\,\mathrm{m}$ 时，车辆仍然频繁出界。
今天的排查重点集中在三个互相耦合的问题：

\begin{enumerate}[label=(\arabic*)]
  \item 原版 alpha-RACER 中 \texttt{TRACK\_WIDTH=0.22}、\texttt{scale=2.0}、\texttt{track\_width=1.0}、以及 \texttt{\(\pm 0.44\)} 到底分别表示什么；
  \item 当前项目中的 world 坐标、Frenet 坐标、raceline deviation、MPC reference frame 是否一致；
  \item 为什么 raw 坐标下按照 $0.22\,\mathrm{m}$ 半宽控制时，MPC 精度仍然不够。
\end{enumerate}

\section{关键代码证据}

\subsection{Raceline 生成宽度}

原项目 raceline 生成处定义：

\begin{lstlisting}[language=Python]
TRACK_WIDTH = 0.22
\end{lstlisting}

对应代码：
\href{https://github.com/sastry-group/alpha-RACER/blob/main/ref_trajs/calculate_raceline.py#L12}{alpha-RACER/ref\_trajs/calculate\_raceline.py:L12}。

该值被写入 raceline 优化器所需的左右边界宽度：
\href{https://github.com/sastry-group/alpha-RACER/blob/main/ref_trajs/calculate_raceline.py#L30-L31}{calculate\_raceline.py:L30--L31}。

本项目读取的 \texttt{berlin\_2018\_raceline\_with\_speeds.csv} 中，最后一列 \texttt{dev\_m} 的最大绝对值约为 $0.206\,\mathrm{m}$。
这与 raw 坐标下半宽 $0.22\,\mathrm{m}$ 基本吻合，说明 raceline 本身几乎贴着边界。

\subsection{\texttt{params-num.yaml} 中的 scale}

原项目数值仿真配置：

\begin{lstlisting}[language=yaml]
track_info:
  track_width: 1.0
  scale: 2.0

timescale: 2.0
\end{lstlisting}

对应代码：
\href{https://github.com/sastry-group/alpha-RACER/blob/main/simulators/params-num.yaml#L1-L10}{params-num.yaml:L1--L10}，
\href{https://github.com/sastry-group/alpha-RACER/blob/main/simulators/params-num.yaml#L36}{params-num.yaml:L36}。

\subsection{WaypointGenerator 中 scale 的实际作用}

原项目 \texttt{WaypointGenerator} 定义在：
\href{https://github.com/sastry-group/alpha-RACER/blob/main/car_dynamics/car_dynamics/controllers_jax/waypoint.py#L69}{car\_dynamics/controllers\_jax/waypoint.py:L69}。

关键逻辑为：

\begin{lstlisting}[language=Python]
self.raceline = np.array(df_raceline.iloc[:,:4]) + np.array([0, ox, oy, 0])
self.raceline_dev = 2*np.array(df_raceline.iloc[:,4])
self.raceline = self.raceline[:,1:]
self.raceline[:,:2] *= float(yaml_content['track_info']['scale'])

if waypoint_type.find('num') != -1:
    self.path = np.array(df.iloc[:-1,:])*yaml_content['track_info']['scale'] + np.array([0, ox, oy, 0])
else:
    self.path = np.array(df.iloc[:,:]) + np.array([0, ox, oy, 0])
self.path[:,-1] = speed
\end{lstlisting}

对应代码：
\href{https://github.com/sastry-group/alpha-RACER/blob/main/car_dynamics/car_dynamics/controllers_jax/waypoint.py#L129-L139}{waypoint.py:L129--L139}。

由此可得：

\begin{itemize}
  \item raceline 的 $x,y$ 会乘以 \texttt{scale=2.0}；
  \item raceline 的 \texttt{dev\_m} 没有通过通用 \texttt{scale} 变量缩放，而是写成 \texttt{2 * dev\_m}；
  \item raceline 的速度列 \texttt{vx\_mps} 没有乘以 2；
  \item centerline/path 的速度列即使暂时被整体乘以 \texttt{scale}，随后也会被 \texttt{self.path[:,-1] = speed} 覆盖。
\end{itemize}

因此原版中 scale 后的主要效果是：

\[
x,y,s,e \text{ 的空间尺度约扩大 } 2 \text{ 倍，而速度不随之扩大。}
\]

\subsection{\texttt{\(\pm 0.44\)} 的来源}

原项目 \texttt{multi\_car\_comp\_blocking.py} 中有：

\begin{lstlisting}[language=Python]
e_ = -self.raceline_dev[lookahead_idx]
if e_ + shift > 0.44:
    shift = 0.44 - e_
if e_ + shift < -0.44:
    shift = -0.44 - e_
\end{lstlisting}

对应代码：
\href{https://github.com/sastry-group/alpha-RACER/blob/main/car_ros2/car_ros2/multi_car_comp_blocking.py#L496-L500}{multi\_car\_comp\_blocking.py:L496--L500}。

因为 \texttt{raceline\_dev = 2 * dev\_m}，而 raw \texttt{dev\_m} 最大约为 $0.206$，所以 scale 后 deviation 最大约为：

\[
2 \times 0.206 \approx 0.412.
\]

这与 clip 半宽 $0.44$ 基本一致。
因此 \texttt{\(\pm 0.44\)} 更像是 scale 后运行坐标系中的 lateral clip，而不是 raw CSV 坐标下的半宽。

\subsection{\texttt{track\_width: 1.0} 的用途}

\texttt{WaypointGenerator} 中左右边界由以下代码生成：

\begin{lstlisting}[language=Python]
self.left_boundary = self.calc_shifted_traj(
    self.path[:,1:],
    yaml_content['track_info']['track_width']/2. + 0.15
)
self.right_boundary = self.calc_shifted_traj(
    self.path[:,1:],
    -yaml_content['track_info']['track_width']/2. - 0.2
)
\end{lstlisting}

对应代码：
\href{https://github.com/sastry-group/alpha-RACER/blob/main/car_dynamics/car_dynamics/controllers_jax/waypoint.py#L143-L144}{waypoint.py:L143--L144}。

进一步搜索发现，\texttt{left\_boundary/right\_boundary} 主要用于 ROS Path 发布和绘图，例如：
\href{https://github.com/sastry-group/alpha-RACER/blob/main/car_ros2/car_ros2/multi_car_comp_blocking.py#L1471-L1488}{multi\_car\_comp\_blocking.py:L1471--L1488}。

因此 \texttt{track\_width: 1.0} 更像可视化或场景边界参数，不是 raceline optimizer 的严格半宽，也不是 policy 中 \texttt{\(\pm 0.44\)} 的直接来源。

\section{当前项目中的坐标标准}

本项目仓库：
\href{https://github.com/CrazalDaver/racer-LLM}{CrazalDaver/racer-LLM}。
注意：若本地改动尚未 push，GitHub 行号可能不能完全反映当前工作区。

\subsection{World 坐标}

车辆状态定义在：
\href{https://github.com/CrazalDaver/racer-LLM/blob/main/state.py}{state.py}。

\begin{lstlisting}[language=Python]
CarState(x, y, vx, vy, psi, omega)
\end{lstlisting}

其中：

\begin{center}
\begin{tabular}{lll}
\toprule
变量 & 坐标系 & 单位 \\
\midrule
\texttt{x, y} & world frame & m \\
\texttt{psi} & world heading & rad \\
\texttt{vx, vy} & vehicle/body velocity & m/s \\
\texttt{omega} & yaw rate & rad/s \\
\bottomrule
\end{tabular}
\end{center}

\subsection{Track Frenet 坐标}

本项目几何投影实现在：
\href{https://github.com/CrazalDaver/racer-LLM/blob/main/track.py}{track.py}。

核心定义为：

\[
p_y = (p_{\mathrm{world}} - p_{\mathrm{center}})^\top n,
\quad
n = (-\sin\psi_{\mathrm{center}}, \cos\psi_{\mathrm{center}}).
\]

即：

\begin{itemize}
  \item \texttt{px/s}: 沿 centerline 的累计弧长，单位 m；
  \item \texttt{py/e}: 相对 centerline 法向偏移，单位 m；
  \item \texttt{py > 0}: 位于 centerline 前进方向左侧；
  \item \texttt{py < 0}: 位于 centerline 前进方向右侧。
\end{itemize}

\subsection{Raceline deviation}

实验确认：

\[
\mathrm{project}(\mathrm{raceline\_x}, \mathrm{raceline\_y}).p_y \approx -\mathrm{dev\_m}.
\]

因此当前代码中使用：

\begin{lstlisting}[language=Python]
p_y = -point["dev_m"]
\end{lstlisting}

从符号上是合理的。
相关逻辑在：
\href{https://github.com/CrazalDaver/racer-LLM/blob/main/policies.py}{policies.py}。

\section{今天发现的关键问题}

\subsection{raw 坐标下速度过激}

原版数值仿真中，\texttt{scale=2.0} 使轨道空间尺度扩大约 2 倍，而 raceline 速度没有随之扩大。
这意味着：

\[
\kappa_{\mathrm{scaled}} \approx \frac{1}{2}\kappa_{\mathrm{raw}}.
\]

在相同速度 $v$ 下，横向加速度需求：

\[
a_y = v^2 \kappa.
\]

因此 scale 后同样的 $5\,\mathrm{m/s}$ 会更容易过弯。
而当前项目如果直接使用 raw CSV 坐标，不做 scale，同时仍使用 raceline 中接近 $5\,\mathrm{m/s}$ 的速度，则弯道横向加速度需求约为原版 scale 后的 2 倍。

这很可能是当前 $0.22\,\mathrm{m}$ 半宽下屡屡出界的核心原因之一。

\subsection{raceline 几乎贴边}

raw 坐标下：

\[
\max|\mathrm{dev\_m}| \approx 0.206,
\quad
w_{\mathrm{half}} = 0.22.
\]

安全余量只有：

\[
0.22 - 0.206 = 0.014\,\mathrm{m}.
\]

这只有约 1.4 cm。即使 MPC 和投影完全正确，DBM 模型误差、离散时间、转向限制、速度过高都会让车辆越界。

\subsection{MPC 模型与 DBM 环境不一致}

当前 MPC 使用 kinematic bicycle 预测：
\href{https://github.com/CrazalDaver/racer-LLM/blob/main/mpc_controller.py}{mpc\_controller.py}。

而 DBM 环境状态为：

\[
[x,y,\psi,v_x,v_y,\omega].
\]

对应代码：
\href{https://github.com/CrazalDaver/racer-LLM/blob/main/vehicle_model.py}{vehicle\_model.py}。

因此在 DBM 下，即使 kinematic MPC 预测轨迹在边界内，真实 roll-out 也可能由于侧滑和 yaw dynamics 偏离 raceline。

\subsection{投影窗口导致的 py 虚高}

今天还发现一个独立问题：旧版车辆投影逻辑会在上一帧 \texttt{progress\_s} 附近局部搜索。
在 DBM 出现回退或靠近其他 track branch 时，记录的 \texttt{py} 可能不是全局最近点 deviation。

最新一次 \texttt{simulation\_trace.npz} 中曾出现：

\begin{center}
\begin{tabular}{llll}
\toprule
时间 & 记录的 \texttt{px, py} & 全局最近 \texttt{s, e} & 结论 \\
\midrule
$23.92\,\mathrm{s}$ & $(21.76, 6.65)$ & $(13.16, 0.053)$ & \texttt{py=6.65} 是局部投影滞后导致 \\
\bottomrule
\end{tabular}
\end{center}

随后车辆投影逻辑已改为：

\begin{lstlisting}[language=Python]
s_wrapped, e = track.project(state.x, state.y)
return s_hint + progress_delta(s_wrapped, s_hint, track.total_length), e
\end{lstlisting}

即：

\begin{itemize}
  \item \texttt{py/e}: 来自全局最近点；
  \item \texttt{px/s}: 使用上一帧进度做 unwrap，避免跨起终点时跳变。
\end{itemize}

相关代码：
\href{https://github.com/CrazalDaver/racer-LLM/blob/main/simulation.py}{simulation.py}。

\section{目前结论}

\begin{enumerate}[label=(\arabic*)]
  \item \texttt{calculate\_raceline.py} 中的 \texttt{TRACK\_WIDTH=0.22} 是 raw CSV 坐标下 raceline 优化半宽。
  \item \texttt{params-num.yaml} 中的 \texttt{scale=2.0} 会把运行时空间坐标放大 2 倍。
  \item 原项目中的 \texttt{\(\pm 0.44\)} 与 $0.22 \times 2$ 对应，属于 scale 后运行坐标系，不应直接当作 raw 坐标下的半宽。
  \item \texttt{track\_width: 1.0} 主要用于边界可视化或场景 Path 发布，不是 raceline optimizer 的严格半宽。
  \item 原版 scale 后速度并没有乘 2，因此同样速度在更大轨道上更容易过弯。
  \item 当前项目若不做 scale，却直接使用 raw 几何和原 raceline 速度，等价于在更小、更弯、更窄的赛道上用相同速度跑，出界风险显著增加。
\end{enumerate}

\section{建议的下一步}

为了让最简单 MPC 控制先稳定跟随 raceline，建议按以下顺序排查：

\begin{enumerate}[label=(\arabic*)]
  \item 保持 raw 几何半宽：
  \[
  w_{\mathrm{half}} = 0.22.
  \]
  \item 参考轨迹内缩，而不是贴着 raceline 最大 deviation 跑：
  \[
  |p_y^{\mathrm{ref}}| \leq 0.17 \sim 0.19.
  \]
  \item 将速度上限先降到：
  \[
  v_{\max} = 3.0 \sim 4.0\,\mathrm{m/s}.
  \]
  \item 先用 kinematic 环境确认 MPC 能稳定不出界，再切 DBM。
  \item 若要复现原版数值环境，应考虑整体引入 \texttt{scale=2.0}，同时一致缩放 centerline、raceline、deviation、车辆位置和动力学解释，而不是只改 track width。
\end{enumerate}

\section{一句话摘要}

当前最可能的真相是：

\begin{quote}
原版 alpha-RACER 的数值实验是在 scale 后的更大轨道上用未放大的速度跑；而当前项目直接使用 raw 小轨道和同样速度，导致曲率需求更高、raceline 又几乎贴边，因此 $0.22\,\mathrm{m}$ 半宽下很容易出界。
\end{quote}

\end{document}
